% Source: https://artofproblemsolving.com/community/c3h2247769_300th_post_problems
% Stable identity: aops-topic-2247769-q03
\textbf{(a)} If $A+B=99$ or $A\times B=99$ with positive integers $A,B$, then $A,B\le99$, so together they have at most $4$ digits, not $10$.

\textbf{(b)} From $A=99B$ we get $10B\le A<100B$, so $A$ has one or two more digits than $B$. If $B$ has $m$ digits, the total number of digits is $2m+1$ or $2m+2$; it equals $10$ only for $2m+2=10$, i.e. $B$ has $4$ digits and $A$ has $6$.

The digit sums $S(A)$ and $S(B)$ add up to $0+1+\cdots+9=45$. Since $9\mid99B=A$, also $9\mid S(A)$, and then $9\mid45-S(A)=S(B)$.

\textbf{(c)} $5013\cdot99=501300-5013=496287$, so $496287\div5013=\boxed{99}$. The digits used are $4,9,6,2,8,7$ and $5,0,1,3$: each of $0,\ldots,9$ exactly once. (Here $S(A)=36$ and $S(B)=9$, as (b) predicts.)

(A computer search over all ways of splitting the ten digits into two numbers shows that this is the only solution, and that the sign $-$ is also impossible.)
